\chapter{Hint and Solution}

\begin{problem}
    \Cref{code: n-queen}의 시간 복잡도는 무엇인가?
\end{problem}

\begin{sol}
    사실 백트레킹 문제의 시간복잡도를 논의하기는 어렵다. \Cref{code: n-queen}를 보면 \texttt{isvalid} 함수를 통과하는지 여부에 따라 다음 재귀 호출이 실행될 수도 있고 실행되지 않을 수도 있기 떄문이다. 다만 확실한 것은 재귀 함수의 호출 횟수가 $T(N)$이면 시간복잡도는 $O(N \cdot T(n))$이라는 사실이다.

    고로 $T(n)$의 상한을 잡아 시간복잡도의 상한을 잡자. 대각선을 고려하지 않았을 때, 매 번 선택되는 열이 다르므로 $T(N) \leq N!$을 쉽게 보일 수 있다. 따라서 시간복잡도의 상한은 $O(N \times N!)$이다.

    Cf. 실제로 상한과 실제 실행 시간이 다름을 3번 문제에서 확인할 수 있을 것이다.
\end{sol}